% Ad M. Brutum 1.5 (ad-brutum-01-05) --- English translation
% Translated by Claude (Anthropic) from the Latin (Perseus).
% Style guide: STYLE.md.

\ciceroLetterOpener{Cicero Bruto salutem}

\ciceroSection{1} On the fifth day before the kalends of May, when motions were being delivered about the prosecution of those who have been declared public enemies, Servilius spoke also concerning Ventidius, and to the effect that Cassius should hunt down Dolabella. As for you, the motion was that, if you should judge it useful and consistent with the public interest, you should pursue Dolabella by war; but if you could not do that without disadvantage to the public interest, or did not consider it consistent with the public interest, that you should keep your army in the same regions. The Senate could have paid you no higher compliment than to leave it to your own judgement what would seem to you most conducive to the state. For my part, this is my view: if Dolabella has a force, has a camp, has anywhere at all to make a stand, it bears on your purpose and on your standing to hunt him down.

\ciceroSection{2} Of our friend Cassius's resources we knew nothing. For no letter has come from him personally, nor was anything reported that we could hold for certain. How much it matters, however, that Dolabella be crushed, you surely understand --- both that he should pay the penalty for his crime, and that the brigand-captains may have no rallying-point to which to flock after their flight from Mutina. And that this has been my view for some time you may recall from earlier letters of mine; though at that time the haven for any flight, and the safeguard of our survival, lay in your camp and in your army. So much the more now, freed (as I hope) from peril, we ought to be wholly engaged in the crushing of Dolabella. But these things you will weigh more carefully and decide more wisely; let us know, if it seems good to you, what you have determined and what you are doing.

\ciceroSection{3} I want our Cicero to be co-opted into your college. I judge that, in general, in the priestly comitia, the case of an absent man can be entertained; for in fact it has been done before. Gaius Marius, while he was in Cappadocia, was made an augur under the \textit{Lex Domitia} (a law on the election of priests), and no later law has restricted this. There is also in the \textit{Lex Iulia} --- the most recent law on priesthoods --- this clause: \textsc{whoever shall stand, or whose case shall be considered}. The wording shows openly that the case can be considered even of one who is not present. I have written to him on this matter, telling him to follow your judgement here as in all things; but you will have to decide about Domitius and about our friend Cato. Yet although it is permissible for the case of an absent man to be considered, still everything goes more easily for those present. And if you should decide that you must go to Asia, there will be no way of summoning our young people back for the elections.

\ciceroSection{4} On the whole, while Pansa was alive we expected everything to move more quickly: for he would at once have got himself a colleague substituted, and then the priestly comitia would have come before those for the praetorship. As it is, I foresee a long delay because of the auspices. As long as there is only one patrician magistrate, the auspices cannot return to the patres. The disturbance is serious indeed. I should like you to inform me what you think on the whole matter. The third day before the nones of May.
